\documentclass[11pt]{article}
\usepackage[margin=1.1in]{geometry}
\usepackage{amsmath,amssymb,amsthm,mathtools,booktabs,array,tabularx,longtable}
\usepackage[dvipsnames]{xcolor}
\usepackage[colorlinks=true,linkcolor=MidnightBlue,citecolor=MidnightBlue,urlcolor=MidnightBlue]{hyperref}
\usepackage{fancyhdr}
\pagestyle{fancy}\fancyhf{}\renewcommand{\headrulewidth}{0pt}
\fancyfoot[R]{\thepage}
\fancyfoot[L]{\scriptsize\color{gray}Generated by the Transport Lab from approved parts of its store; not a deposited paper. Papers of record: \href{https://zenodo.org/communities/transport-papers}{Zenodo, Transport Papers}}
\fancypagestyle{plain}{\fancyhf{}\fancyfoot[R]{\thepage}\fancyfoot[L]{\scriptsize\color{gray}Generated by the Transport Lab from approved parts of its store; not a deposited paper. Papers of record: \href{https://zenodo.org/communities/transport-papers}{Zenodo, Transport Papers}}\renewcommand{\headrulewidth}{0pt}}
\newtheorem{theorem}{Theorem}[section]
\newtheorem{proposition}[theorem]{Proposition}
\newtheorem{lemma}[theorem]{Lemma}
\newtheorem{corollary}[theorem]{Corollary}
\newtheorem{conjecture}[theorem]{Conjecture}
\newtheorem{refutation}[theorem]{Refutation}
\newtheorem{prediction}[theorem]{Prediction}
\newtheorem{principle}[theorem]{Principle}
\newtheorem{observation}[theorem]{Observation}
\newtheorem{axiom}[theorem]{Axiom}
\theoremstyle{definition}
\newtheorem{definition}[theorem]{Definition}
\newtheorem{fact}[theorem]{Fact}
\newtheorem{assembled}[theorem]{Assembly}
\newtheorem{postulate}[theorem]{Postulate}
\newtheorem{premise}[theorem]{Premise}
\newtheorem{hypothesis}[theorem]{Hypothesis}
\newtheorem{remark}[theorem]{Remark}
\newtheorem{example}[theorem]{Example}
\theoremstyle{plain}
\setlength{\marginparwidth}{0.8in}
\newcommand{\changedmark}[2]{\marginpar{\raggedright\scriptsize\color{#1}#2}}
\providecommand{\Aa}{}\renewcommand{\Aa}{\mathcal{A}}
\providecommand{\Ad}{}\renewcommand{\Ad}{\widetilde{\mathrm{Ad}}}
\providecommand{\Aev}{}\renewcommand{\Aev}{\mathcal{A}^{+}}
\providecommand{\Aod}{}\renewcommand{\Aod}{\mathcal{A}^{-}}
\providecommand{\B}{}\renewcommand{\B}{\mathrm{B}}
\providecommand{\C}{}\renewcommand{\C}{\mathbb{C}}
\providecommand{\CAR}{}\renewcommand{\CAR}{\mathrm{CAR}}
\providecommand{\CS}{}\renewcommand{\CS}{{\mathbb{S}^{+}}}
\providecommand{\CSm}{}\renewcommand{\CSm}{{\mathbb{S}^{-}}}
\providecommand{\Cl}{}\renewcommand{\Cl}{\mathrm{Cl}}
\providecommand{\Con}{}\renewcommand{\Con}{\mathcal{C}}
\providecommand{\Dl}{}\renewcommand{\Dl}{\Delta}
\providecommand{\EN}{}\renewcommand{\EN}{E_{N}}
\providecommand{\End}{}\renewcommand{\End}{\mathrm{End}}
\providecommand{\F}{}\renewcommand{\F}{\mathbb{F}}
\providecommand{\G}{}\renewcommand{\G}{\Gamma}
\providecommand{\HH}{}\renewcommand{\HH}{\mathbb{H}}
\providecommand{\IS}{}\renewcommand{\IS}{I\sigma_3}
\providecommand{\J}{}\renewcommand{\J}{\mathbf{J}}
\providecommand{\LH}{}\renewcommand{\LH}{L_{H}}
\providecommand{\Mat}{}\renewcommand{\Mat}{\mathrm{M}}
\providecommand{\Mo}{}\renewcommand{\Mo}{M^{\circ}}
\providecommand{\Mp}{}\renewcommand{\Mp}{M^{+}}
\providecommand{\Msun}{}\renewcommand{\Msun}{M_{\odot}}
\providecommand{\Mt}{}\renewcommand{\Mt}{\widetilde{M}}
\providecommand{\Ncone}{}\renewcommand{\Ncone}{\mathcal{N}}
\providecommand{\Neff}{}\renewcommand{\Neff}{N_{\mathrm{eff}}}
\providecommand{\PCS}{}\renewcommand{\PCS}{\mathbb{P}{\mathbb{S}^{+}}}
\providecommand{\PT}{}\renewcommand{\PT}{\mathbb{PT}}
\providecommand{\Pin}{}\renewcommand{\Pin}{\mathrm{Pin}}
\providecommand{\Q}{}\renewcommand{\Q}{\mathbb{Q}}
\providecommand{\R}{}\renewcommand{\R}{\mathbb{R}}
\providecommand{\Ric}{}\renewcommand{\Ric}{\mathrm{Ric}}
\providecommand{\SBH}{}\renewcommand{\SBH}{S_{\mathrm{BH}}}
\providecommand{\Sc}{}\renewcommand{\Sc}{\mathbb{S}}
\providecommand{\Sch}{}\renewcommand{\Sch}{\mathrm{P}}
\providecommand{\Sig}{}\renewcommand{\Sig}{\Sigma}
\providecommand{\Som}{}\renewcommand{\Som}{S_\omega}
\providecommand{\Sp}{}\renewcommand{\Sp}{\mathbb{S}}
\providecommand{\Sph}{}\renewcommand{\Sph}{\mathrm{S}}
\providecommand{\Spi}{}\renewcommand{\Spi}{S_\pi}
\providecommand{\Spin}{}\renewcommand{\Spin}{\mathrm{Spin}}
\providecommand{\Sth}{}\renewcommand{\Sth}{S_{\mathrm{th}}}
\providecommand{\T}{}\renewcommand{\T}{\mathbb{T}}
\providecommand{\TH}{}\renewcommand{\TH}{T_{\mathrm{H}}}
\providecommand{\Tr}{}\renewcommand{\Tr}{\operatorname{Tr}}
\providecommand{\UHF}{}\renewcommand{\UHF}{\mathrm{UHF}}
\providecommand{\Z}{}\renewcommand{\Z}{\mathbb{Z}}
\providecommand{\Zc}{}\renewcommand{\Zc}{\mathrm{Z}}
\providecommand{\cmg}{}\renewcommand{\cmg}{\mathrm{cm^{2}\,g^{-1}}}
\providecommand{\dS}{}\renewcommand{\dS}{\mathrm{dS}}
\providecommand{\dsc}{}\renewcommand{\dsc}{\mathrm{dsc}}
\providecommand{\dtPG}{}\renewcommand{\dtPG}{dt}
\providecommand{\dtS}{}\renewcommand{\dtS}{dt_{\mathrm{S}}}
\providecommand{\ellv}{}\renewcommand{\ellv}{\ell}
\providecommand{\er}{}\renewcommand{\er}{e_r}
\providecommand{\et}{}\renewcommand{\et}{e_t}
\providecommand{\ev}{}\renewcommand{\ev}{\mathrm{ev}}
\providecommand{\gl}{}\renewcommand{\gl}{\mathfrak{gl}}
\providecommand{\hatV}{}\renewcommand{\hatV}{\widehat{V}}
\providecommand{\hb}{}\renewcommand{\hb}{\bar{h}}
\providecommand{\hf}{}\renewcommand{\hf}{\hat{f}}
\providecommand{\id}{}\renewcommand{\id}{\mathrm{id}}
\providecommand{\kpc}{}\renewcommand{\kpc}{\mathrm{kpc}}
\providecommand{\lP}{}\renewcommand{\lP}{\ell_{P}}
\providecommand{\mD}{}\renewcommand{\mD}{\mathbf{m}_{D}}
\providecommand{\mb}{}\renewcommand{\mb}{\bar m}
\providecommand{\mnu}{}\renewcommand{\mnu}{\Sigma m_{\nu}}
\providecommand{\nInf}{}\renewcommand{\nInf}{n_\infty}
\providecommand{\nO}{}\renewcommand{\nO}{n_0}
\providecommand{\ninf}{}\renewcommand{\ninf}{n_{\infty}}
\providecommand{\nz}{}\renewcommand{\nz}{n_{0}}
\providecommand{\proj}{}\renewcommand{\proj}{\operatorname{proj}}
\providecommand{\pv}{}\renewcommand{\pv}{\pi}
\providecommand{\rad}{}\renewcommand{\rad}{\mathrm{rad}}
\providecommand{\sgn}{}\renewcommand{\sgn}{\operatorname{sgn}}
\providecommand{\shA}{}\renewcommand{\shA}{\mathcal{A}}
\providecommand{\shC}{}\renewcommand{\shC}{\mathcal{C}}
\providecommand{\shE}{}\renewcommand{\shE}{\mathcal{E}}
\providecommand{\shJ}{}\renewcommand{\shJ}{\mathcal{J}}
\providecommand{\shK}{}\renewcommand{\shK}{\mathcal{K}}
\providecommand{\shL}{}\renewcommand{\shL}{L}
\providecommand{\shW}{}\renewcommand{\shW}{\mathcal{W}}
\providecommand{\sll}{}\renewcommand{\sll}{\mathfrak{sl}}
\providecommand{\so}{}\renewcommand{\so}{\mathfrak{so}}
\providecommand{\su}{}\renewcommand{\su}{\mathfrak{su}}
\providecommand{\supp}{}\renewcommand{\supp}{\operatorname{supp}}
\providecommand{\tr}{}\renewcommand{\tr}{\operatorname{tr}}
\providecommand{\uu}{}\renewcommand{\uu}{\mathfrak{u}}
\providecommand{\br}{}\renewcommand{\br}[1]{\langle #1\rangle}
\providecommand{\bra}{}\renewcommand{\bra}[1]{\langle #1|}
\providecommand{\file}{}\renewcommand{\file}[1]{\texttt{#1}}
\providecommand{\hp}{}\renewcommand{\hp}[2]{( #1 \,|\, #2 )}
\providecommand{\ip}{}\renewcommand{\ip}[2]{\langle #1, #2\rangle}
\providecommand{\ket}{}\renewcommand{\ket}[1]{|#1\rangle}
\providecommand{\rel}{}\renewcommand{\rel}[2]{D\bigl(#1 \,\|\, #2\bigr)}
\providecommand{\rev}{}\renewcommand{\rev}[1]{\widetilde{#1}}
\providecommand{\slashed}{}\renewcommand{\slashed}[1]{{#1}\mkern-10mu/}
\title{\bfseries Spacetime in the conformal completion}
\author{Leon Butler\\ \normalsize\href{mailto:lbutler.math@gmail.com}{\texttt{lbutler.math@gmail.com}}\\ \normalsize ORCID: \href{https://orcid.org/0009-0006-4449-5617}{0009-0006-4449-5617}}
\date{Generated by the Transport Lab, 2 November 2026}
\begin{document}
\maketitle
\begin{abstract}
\normalsize
In the conformal completion, two extra directions are added to the algebra of spacetime, and every point becomes a null vector, one whose square is zero. This trail shows that the same algebra already contains the conformal spinor space, which Penrose called twistor space.

The null directions of the completion match the lines of the projective spinor space, the classical Klein correspondence obtained here from the algebra. Penrose's incidence relation, the rule tying points to spinors, turns out to be the completion's own coordinate chart.

Conformal spinors are vacuum states of fermionic ladder pairs, not the operators that act on them. The walk ends with Einstein's vacuum equations read in three equivalent languages.
\end{abstract}
\newpage
\noindent \textbf{Authorship and disclosure. }The founding idea of the series is the author's: that division by zero in a Clifford algebra should be read not as a prohibition but as a transport into a higher-dimensional algebra, and that the transported division need not be accompanied by an inverse operation. The direction of investigation and the critical review of the results are likewise the author's. The mathematical development and the drafting of the text were carried out in extended collaboration with Claude Fable 5 (Anthropic), an AI system.

\noindent\fbox{\parbox{\dimexpr\linewidth-2\fboxsep-2\fboxrule}{\textbf{This is not a deposited paper.} It was generated by the Transport Lab on 2 November 2026 and assembled from approved parts of its store. It follows results across several papers of the series; each is stated with its proof and credited to the paper that first proved it. The official versions of those papers are on Zenodo, in the Transport Papers community, \url{https://zenodo.org/communities/transport-papers}.}}
\newpage

\section*{Introduction}
The walk starts from points as null vectors, builds the space the completed algebra acts on, and then reads the spinor dictionary off it one step at a time. Results a step rests on but that are not on the trail are listed at the end.

\section{Spacetime in the conformal completion}

These results place spacetime, its light cones and its conformal spinors inside a single algebra.

The completion adds a null pair $\nO$ and $\nInf$, the origin and the point at infinity, to the directions of spacetime. A point $x$ is then represented by the null vector \[ X(x) = \nO + x + \tfrac12 Q(x)\,\nInf \]

The product of two such vectors is $X(x)\cdot X(y) = -\tfrac12 Q(x - y)$, so separations are read off the algebra directly. The rotors $R_u = 1 + vu$ that fix a null vector $v$ add like vectors.

\begin{assembled}[Imported: the conformal model and the null rotations]\label{c:angle.imported-conformal-model-null-rotations}
Adding two extra directions to the algebra of space lets every point be written as a null vector, one whose square is zero, and the squared distance between two points can then be read off from a simple product of their vectors. In the same setting, each null direction has a family of rotations that leave it fixed, and these rotations combine exactly like ordinary addition of vectors. Both facts are carried over from earlier work in the series.
\emph{First proved in} T5 (angle note), Fact 1.2 (\href{https://doi.org/10.5281/zenodo.22721790}{doi:10.5281/zenodo.22721790}).
\end{assembled}
\noindent\emph{Rests on:} \ref{c:transport.group-null-rotations}, \ref{c:transport.totality}.

To compute with the completed algebra we need a concrete space for it to act on.

Pick a maximal set $W$ of mutually null complex directions, with $\dim_\C W = 3$, and let vectors act on the exterior algebra $\Lambda^\bullet W$ by wedging and contracting: \[ c(w) = \sqrt2\; w \wedge (\cdot), \qquad c(w') = \sqrt2\; \iota\_{w'} \] This is the only irreducible module of the complexified algebra, the Fock model.

Vectors act oddly, so the module splits into even and odd halves of complex dimension four, the two chiralities $\CS$ and $\CSm$.

\begin{theorem}[The Fock model of the complexified ascended algebra]\label{c:correspondence.fock-model-complexified-ascended-algebra}
The completed algebra of six directions can be represented concretely. Pick three complex directions that are mutually lightlike, and let the algebra act on all combinations of them by adding or removing a direction, much as one creates and destroys particles. This gives an eight-dimensional space on which the algebra acts faithfully and irreducibly, the only such space up to equivalence, and it splits into two four-dimensional halves, the two chiralities.
\emph{First proved in} T9, Thm 3.1 (\href{https://doi.org/10.5281/zenodo.22722290}{doi:10.5281/zenodo.22722290}).
\end{theorem}
\noindent\emph{Rests on:} \ref{c:std-algebra.existence-uniqueness}, \ref{c:std-algebra.complex-clifford-algebra}.

With the conformal spinor space in hand, the next question is how the symmetries of spacetime act on it.

The bivectors close under commutators and act on $\CS$ by an isomorphism \[ \rho : $\mathfrak{so}$(2,4)\_\C \to $\mathfrak{sl}$(\CS) \] So the spin group acts as all of $SL(\CS)$, and any nonzero conformal spinor can be moved to any other.

\begin{proposition}[The spin representation is $\mathfrak{sl}(\CS)$]\label{c:correspondence.spin-representation}
The conformal symmetries of spacetime, extended to complex numbers, can be realised inside the completed Clifford algebra. Acting on the four-dimensional spinor space of the conformal algebra they give exactly the trace-free linear transformations of that space, no more and no fewer. At the level of groups every volume-preserving linear map of that space is achieved, so any nonzero spinor in it can be moved to any other.
\emph{First proved in} T9, Prop 5.2 (\href{https://doi.org/10.5281/zenodo.22722290}{doi:10.5281/zenodo.22722290}).
\end{proposition}
\noindent\emph{Rests on:} \ref{c:std-algebra.vector-squares}, \ref{c:std-algebra.complex-clifford-algebra}.

Now the two pictures can be matched.

The complexified vectors are exactly the bivectors of the spinor space, \[ V\_\C \cong \Lambda\^{}{2}\CS \] with the quadratic form becoming the wedge square of a bivector. A vector is null exactly when its bivector is simple, that is, when it is a line of the projective spinor space $\PCS$. This is the Klein correspondence, obtained from the completion rather than assumed.

\begin{theorem}[The spinor identification]\label{c:correspondence.spinor-identification}
This result matches the six complex directions of the completed spacetime algebra with the six independent plane elements that can be formed inside the four-dimensional conformal spinor space. Under the match, the lightlike directions are exactly those plane elements that are single planes rather than sums of planes, and in projective terms these are lines of that spinor space. So the lightlike directions of the completion, taken up to sign, are lines in the projective spinor space: the Klein correspondence, obtained from the completion.
\emph{First proved in} T9, Thm 3.4 (\href{https://doi.org/10.5281/zenodo.22722290}{doi:10.5281/zenodo.22722290}).
\end{theorem}
\noindent\emph{Rests on:} \ref{c:correspondence.chevalley-pairing}, \ref{c:correspondence.fock-model-complexified-ascended-algebra}, \ref{c:correspondence.imported-forced-completion-fibre-quadric}, \ref{c:std-algebra.vector-squares}, \ref{c:std-algebra.inverse-of-vector}.

Since points are null vectors and null vectors are lines, each point of spacetime should be a line of spinors.

For a linear map $A$ between the two planes $\Spi$ and $\Som$ of the splitting, the bivector of its graph line expands as \[ P(A) = N\_0 + \varphi(A) + (\det A)\, N\_\infty \] Term by term this is the completion's point $X(m(A))$, with the determinant as its quadratic term, so the incidence relation is the completion's chart.

\begin{theorem}[The incidence relation is the chart]\label{c:correspondence.incidence-relation-chart}
Penrose described a spacetime point by a line in a three-dimensional complex projective space, through his incidence relation. This theorem shows that, written out term by term, the coordinates of that line are exactly the null vector by which the completed algebra represents the point, so the incidence relation is the completion's coordinate chart. The quadratic term of the chart is a determinant, the lines meeting a fixed line at infinity are the points at infinity, and the symmetries fixing that line are those of flat spacetime.
\emph{First proved in} T9, Thm 4.2 (\href{https://doi.org/10.5281/zenodo.22722290}{doi:10.5281/zenodo.22722290}).
\end{theorem}
\noindent\emph{Rests on:} \ref{c:correspondence.chiral-planes-summand-spaces}, \ref{c:correspondence.imported-stabilisers-flat-model}, \ref{c:correspondence.skew-lines-splittings-hyperbolic-planes-one}, \ref{c:correspondence.spinor-identification}.

So far everything is complex, but the spacetime we live in is real.

The completion supplies its own antilinear conjugation, unique up to scale, and from it a Hermitian form on $\CS$ of signature $(2,2)$. Real points are the null lines of that form skew to the infinity line, with incidence relation \[ \omega = I\_6\, \hat x\, \pi \] for Hermitian $\hat x$, so the factor in front is the conformal pseudoscalar $I_6$, not a bare imaginary unit.

\begin{proposition}[Reality: the conjugation, the invariant Hermitian form, and the factor $I_6$]\label{c:correspondence.reality-conjugation-twistor-form-factor}
The spinor space of the conformal algebra, as built here, is a complex space, and real spacetime has to be recovered from it by a reality condition. This result builds that condition from the completion itself: a conjugation operation and, from it, a Hermitian form on the spinor space with two positive and two negative directions. Real spacetime points are exactly the lines of spinors on which this form vanishes, and the complex unit in Penrose's incidence relation, here the conformal pseudoscalar, is forced by this reality condition.
\emph{First proved in} T9, Prop 4.5 (\href{https://doi.org/10.5281/zenodo.22722290}{doi:10.5281/zenodo.22722290}).
\end{proposition}
\noindent\emph{Rests on:} \ref{c:correspondence.chevalley-pairing}, \ref{c:correspondence.incidence-relation-chart}, \ref{c:correspondence.fock-model-complexified-ascended-algebra}, \ref{c:std-algebra.inverse-of-vector}, \ref{c:std-algebra.reflection}.

Having placed conformal spinors in the algebra, we can ask what kind of object they are there.

Each nonzero spinor $u$ has an annihilator \[ W\_u = \{ v \in V\_\C : c(v)u = 0 \} \] which is a maximal null subspace, so every conformal spinor is pure. The spinor is the Fock vacuum of the three ladder pairs built on $W_u$, and every other state is reached from it by those pair operators.

Conformal spinors are therefore states that operators act on, not the operators themselves.

\begin{theorem}[Conformal spinors are the states of three ladder pairs]\label{c:correspondence.twistors-are-states-three-ladder-pairs}
The four-component spinors of the conformal algebra are the objects Penrose built his theory on. This result shows that, inside the completed algebra of the series, each such spinor is the empty state (the vacuum) of a system of three fermionic ladder pairs that the spinor itself singles out. Every other state of the module is reached from it by those ladder operators, so these spinors are the states the operators act on, not the operators themselves.
\emph{First proved in} T9, Thm 5.3 (\href{https://doi.org/10.5281/zenodo.22722290}{doi:10.5281/zenodo.22722290}).
\end{theorem}
\noindent\emph{Rests on:} \ref{c:correspondence.fock-model-complexified-ascended-algebra}, \ref{c:correspondence.imported-ladder-pairs-fock-modules-over}, \ref{c:correspondence.spin-representation}, \ref{c:std-algebra.vector-squares}, \ref{c:std-algebra.blade-basis}.

Putting the last two steps together gives a new picture of spacetime itself.

Every point is null, so its vector acts on the Fock module as a nilpotent: \[ c(X(x))\^{}{2} = Q(X(x)) = 0 \] Its kernel is a line of vacua, and two points are null-separated exactly when their lines share a vacuum.

\begin{proposition}[Spacetime as the manifold of vacuum lines]\label{c:one-relation.spacetime-manifold-vacuum-lines}
In this picture every point of spacetime is a light-like vector, so in the algebra each point behaves like a zero divisor. Each point also carries a whole line of possible vacuum states, and two points are light-like separated exactly when their lines share a vacuum. The result then reads gravity as how these vacua change from point to point and quantum mechanics as the excitations above a chosen vacuum, with singularities where the family of vacua breaks down.
\emph{First proved in} T16, Prop 4.1 (\href{https://doi.org/10.5281/zenodo.22723212}{doi:10.5281/zenodo.22723212}).
\end{proposition}
\noindent\emph{Rests on:} \ref{c:one-relation.imported-pure-spinors-incidence-chart}, \ref{c:one-relation.imported-witts-extension-theorem}, \ref{c:one-relation.two-presentations}, \ref{c:sheaf.localised-transport}, \ref{c:sheaf.one-sided-completion}, \ref{c:sheaf.one-sided-completion-singularity}, \ref{c:std-algebra.inverse-of-vector}.

For the vacua to vary from point to point we need a rule for carrying spinors along.

Assuming the conformal manifold is spin, the spinor connection induced from the tractor gauge field has exactly the form of Penrose's local transport of four-component spinors. The two agree once one imported sign and normalisation convention for the Schouten term is matched.

\begin{proposition}[The conformal spinor connection is Penrose's local twistor transport]\label{c:correspondence.local-twistor-transport-ascended-spinor-connection}
Two older formalisms describe the geometry of conformal spacetime: tractors, a six-component bundle with its own connection, and Penrose's four-component bundle with its own transport rule. Assuming the conformal spacetime is spin, so that the tractor bundle has spinors (locally they always exist), this result shows that those spinors, built from the completion's Fock model, carry a connection with exactly the shape of Penrose's transport rule. The two agree completely once one sign and normalisation convention for the curvature term is matched, and that convention is the only thing imported.
\emph{First proved in} T9, Prop 4.10 (\href{https://doi.org/10.5281/zenodo.22722290}{doi:10.5281/zenodo.22722290}).
\end{proposition}
\noindent\emph{Rests on:} \ref{c:correspondence.imported-tractor-connection-local-twistor-transport-external}.

With transport in place, Einstein's equations can be stated in any of the three languages met so far.

On the region where the gravitational field $(\hb, \Omega)$ on flat spacetime is invertible, the vacuum and Einstein equations say that the infinity bivector \[ I\_{\alpha\beta} = \Psi(I\_g) \in \Lambda\^{}{2}\CS \] is covariantly constant. It is simple, a genuine line of $\PCS$, exactly when the cosmological constant vanishes, by the same rule that matched null vectors with lines.

\begin{corollary}[One statement in three languages]\label{c:correspondence.one-statement-three-languages}
In this series gravity is a pair of fields on flat spacetime, and Einstein's vacuum equations can be phrased as the statement that one vector of the six-dimensional conformal model, the scale tractor, is constant as those fields carry it around. Translating that object into the spinors of the conformal algebra gives a bivector of that spinor space, the infinity bivector. This result shows that the three formulations say the same thing, and that the familiar fact that this bivector is a genuine line exactly when the cosmological constant vanishes is the same rule that matches lightlike vectors with lines.
\emph{First proved in} T9, Cor 4.8 (\href{https://doi.org/10.5281/zenodo.22722290}{doi:10.5281/zenodo.22722290}).
\end{corollary}
\noindent\emph{Rests on:} \ref{c:correspondence.imported-vacuum-equations-scale-tractor}, \ref{c:correspondence.spinor-identification}.

\section{Results used from elsewhere}
The proofs above use the following results of the series, stated briefly; each is proved in the paper named.

\begin{fact}[The Chevalley pairing]\label{c:correspondence.chevalley-pairing}
The Chevalley pairing on the Fock module pairs the two chiralities perfectly, is antisymmetric, makes every vector self-transpose and is infinitesimally $\Spin$-invariant.
\emph{Source:} T9, Lem 3.3 (\href{https://doi.org/10.5281/zenodo.22722290}{doi:10.5281/zenodo.22722290}).
\end{fact}

\begin{fact}[The chiral planes, and the summand spaces]\label{c:correspondence.chiral-planes-summand-spaces}
The two planes $\Som$, $\Spi$ of the conformal spinor space are the chiral doublets of the complexified Lorentz algebra, their tensor product is complexified Minkowski space, and each equivariant map from a direction and one plane to the other is unique up to scale.
\emph{Source:} T9, Lem 4.4 (\href{https://doi.org/10.5281/zenodo.22722290}{doi:10.5281/zenodo.22722290}).
\end{fact}

\begin{fact}[Imported: the forced completion and its fibre quadric]\label{c:correspondence.imported-forced-completion-fibre-quadric}
Division by a type-N zero divisor in $\Cl(1,3)$ has no element-level resolution, its unique minimal resolving completion is $\Cl(1,3) \hookrightarrow \Cl(2,4)$, and the resolving assignments lie on the null quadric of $\R^{2,4}$.
\emph{Source:} T9, Fact 2.1 (\href{https://doi.org/10.5281/zenodo.22722290}{doi:10.5281/zenodo.22722290}).
\end{fact}

\begin{fact}[Imported: ladder pairs and Fock modules over the reals]\label{c:correspondence.imported-ladder-pairs-fock-modules-over}
Each hyperbolic pair in a real Clifford algebra gives a fermionic ladder pair, and the ladder pairs of a Witt decomposition act as a Fock module on the exterior algebra of the creation span, irreducibly when the form is split.
\emph{Source:} cite[ S7] ButlerT1 ;  cite[ S 4] ButlerT3.
\end{fact}

\begin{fact}[Imported: stabilisers in the flat model]\label{c:correspondence.imported-stabilisers-flat-model}
In $\mathfrak{so}(2,4)$ the stabiliser of the null vector at infinity is the Poincar\'e algebra, the stabiliser of its ray is the eleven-dimensional maximal parabolic, and its nilradical is the translations, which are the no-go theorem's transvections.
\emph{Source:} T9, Fact 4.3 (\href{https://doi.org/10.5281/zenodo.22722290}{doi:10.5281/zenodo.22722290}).
\end{fact}

\begin{fact}[Imported: the tractor connection, and the conformal spinor connection (local twistor transport), external ingredients]\label{c:correspondence.imported-tractor-connection-local-twistor-transport-external}
The literature sources for the tractor connection and the conformal spinor connection: Penrose and Rindler, Bailey, Eastwood and Gover, Curry and Gover, Penrose and MacCallum, and Dighton.
\emph{Source:} R. Penrose and W. Rindler, Spinors and Space-Time, Vols. 1--2, Cambridge University Press (1984, 1986)., Vol.~2, Ch.~6; T. N. Bailey, M. G. Eastwood and A. R. Gover, Thomas's structure bundle for conformal, projective and related structures, Rocky Mountain J. Math. 24 (1994), 1191--1217.; S. N. Curry and A. R. Gover, An introduction to conformal geometry and tractor calculus, with a view towards analysis on manifolds, in Asymptotic Analysis in General Relativity, London Mathematical Society Lecture Note Series 443, Cambridge University Press (2018), 86--170.; R. Penrose and M. A. H. MacCallum, Twistor theory: an approach to the quantisation of fields and space-time, Phys. Rep. 6 (1973), 241--316.; K. Dighton, An introduction to the theory of local twistors, Int. J. Theor. Phys. 11 (1974), 31--43..
\end{fact}

\begin{fact}[Imported: the vacuum equations and the scale tractor]\label{c:correspondence.imported-vacuum-equations-scale-tractor}
On the invertible region of a gravitational field on flat spacetime the vacuum equations are the parallelism of the scale tractor for the tractor gauge field, and a parallel section $(1, 0, \rho)$ is exactly an Einstein scale with $\Lambda = 6\rho$ and squared norm $\Lambda/3$.
\emph{Source:} T9, Fact 4.7 (\href{https://doi.org/10.5281/zenodo.22722290}{doi:10.5281/zenodo.22722290}).
\end{fact}

\begin{fact}[Skew lines, splittings, hyperbolic planes: one datum]\label{c:correspondence.skew-lines-splittings-hyperbolic-planes-one}
An ordered pair of skew lines in the projective spinor space, a normalised splitting of the conformal spinor space into two planes and a normalised null pair in $V_\C$ determine one another, with the mixed slot equal to the orthogonal complement of the null pair.
\emph{Source:} T9, Prop 4.1 (\href{https://doi.org/10.5281/zenodo.22722290}{doi:10.5281/zenodo.22722290}).
\end{fact}

\begin{fact}[Imported: pure spinors and the incidence chart]\label{c:one-relation.imported-pure-spinors-incidence-chart}
Imported: every conformal spinor is a pure spinor, the Fock vacuum of its own three-pair system, and spacetime points are null vectors of $\C^6$, equivalently lines of the projective spinor space.
\emph{Source:} T16, Fact 2.3 (\href{https://doi.org/10.5281/zenodo.22723212}{doi:10.5281/zenodo.22723212}).
\end{fact}

\begin{fact}[Imported: Witt's extension theorem]\label{c:one-relation.imported-witts-extension-theorem}
Over a field of characteristic not two, every isometry between subspaces of a nondegenerate quadratic space extends to the whole space, so the isometry group is transitive on subspaces of each isometry type.
\emph{Source:} E. Witt, Theorie der quadratischen Formen in beliebigen K"orpern, J. Reine Angew. Math. 176 (1937) 31--44.; T. Y. Lam, Introduction to Quadratic Forms over Fields, Graduate Studies in Mathematics 67, American Mathematical Society (2005)., Ch.~I, \S4.
\end{fact}

\begin{fact}[The two presentations]\label{c:one-relation.two-presentations}
Read on an orthogonal basis the Clifford relation is a metric, and read on a Witt basis it is the canonical anticommutation relations, so the geometric and quantum readings are two charts on one algebra.
\emph{Source:} T16, Prop 3.1 (\href{https://doi.org/10.5281/zenodo.22723212}{doi:10.5281/zenodo.22723212}).
\end{fact}

\begin{fact}[Localised transport]\label{c:sheaf.localised-transport}
The degenerate Clifford sheaf embeds canonically in the completed sheaf, and on the wall the radical covector acquires a partner with which it forms a ladder pair that exists exactly over the wall.
\emph{Source:} T4, Cor 4.5 (\href{https://doi.org/10.5281/zenodo.22721669}{doi:10.5281/zenodo.22721669}).
\end{fact}

\begin{fact}[One-sided completion]\label{c:sheaf.one-sided-completion}
Under tameness, with signature $(p,q)$, $p \geq 1$, on the physical side and a trivial radical line bundle, coupling an adjoined hyperbolic plane to the dying direction gives an ambient form of constant signature $(p{+}1, q{+}1)$ on the whole physical side, wall included.
\emph{Source:} T4, Thm 4.4 (\href{https://doi.org/10.5281/zenodo.22721669}{doi:10.5281/zenodo.22721669}).
\end{fact}

\begin{fact}[One-sided completion at the singularity]\label{c:sheaf.one-sided-completion-singularity}
At the blown-up Schwarzschild singularity, coupling an adjoined hyperbolic plane to the dying direction with $\mu\nu < 0$ and bump inside the horizon gives an ambient form of constant signature $(2,4)$ everywhere, boundary included, with no shadow wall and a ladder pair existing exactly over the boundary.
\emph{Source:} T4, Thm 8.12 (\href{https://doi.org/10.5281/zenodo.22721669}{doi:10.5281/zenodo.22721669}).
\end{fact}

\begin{fact}[The blade basis, degenerate forms included]\label{c:std-algebra.blade-basis}
For any basis of $V$ the $2^{n}$ increasing products $e_I$ form a basis of $\Cl(V, Q)$, degenerate forms included; so $V \hookrightarrow \Cl(V, Q)$.
\emph{Source:} the algebra core of the Transport Lab, proved there with an exact check; standard, see C. Chevalley, The Algebraic Theory of Spinors, Columbia University Press (1954), Ch. II (the Clifford algebra of a quadratic space).
\end{fact}

\begin{fact}[The complex Clifford algebra, its classification modulo 2, its Fock model and its pairing]\label{c:std-algebra.complex-clifford-algebra}
$\Cl(V_\C)$ has a blade basis and only depends on $n$: $\Cl_{2k}(\C) \cong \End(\Lambda W) \cong \Mat_{2^{k}}(\C)$ on any maximal totally null $W$, $\Cl_{2k+1}(\C) \cong \Mat_{2^{k}}(\C)^{\oplus 2}$; $\operatorname{top}(\tau x \wedge y)$ is a $\C$-bilinear invariant pairing, antisymmetric for $k = 3$.
\emph{Source:} the algebra core of the Transport Lab, proved there with an exact check; standard, see C. Chevalley, The Algebraic Theory of Spinors, Columbia University Press (1954), Ch. II and III (Clifford algebras over a field; spinors, pure spinors).
\end{fact}

\begin{fact}[The Clifford algebra exists and is unique]\label{c:std-algebra.existence-uniqueness}
$\Cl(V, Q) = T(V)/\langle v \otimes v - Q(v) \rangle$ exists and is unique up to unique isomorphism; the wedge-plus-contraction operators on $\Lambda V$ give a concrete representation.
\emph{Source:} the algebra core of the Transport Lab, proved there with an exact check; standard, see C. Chevalley, The Algebraic Theory of Spinors, Columbia University Press (1954), Ch. II (the Clifford algebra of a quadratic space).
\end{fact}

\begin{fact}[Which vectors and versors are invertible]\label{c:std-algebra.inverse-of-vector}
$v$ is invertible iff $Q(v) \ne 0$, with $v^{-1} = v/Q(v)$; a versor $g = u_1 \cdots u_k$ has $g\tilde{g} = \prod Q(u_i)$.
\emph{Source:} the algebra core of the Transport Lab, proved there with an exact check; standard, see C. Doran and A. Lasenby, Geometric Algebra for Physicists, Cambridge University Press (2003), Ch. 4 (foundations of geometric algebra: rotors, versors).
\end{fact}

\begin{fact}[Reflection in a non-null vector]\label{c:std-algebra.reflection}
For $Q(u) \ne 0$, $x \mapsto -uxu^{-1} = x - 2\frac{B(x,u)}{Q(u)}u$ is the reflection in $u^{\perp}$: an involutive isometry of determinant $-1$ fixing $\operatorname{rad} V$.
\emph{Source:} the algebra core of the Transport Lab, proved there with an exact check; standard, see C. Doran and A. Lasenby, Geometric Algebra for Physicists, Cambridge University Press (2003), Ch. 4 (foundations of geometric algebra: rotors, versors).
\end{fact}

\begin{fact}[The product of two vectors]\label{c:std-algebra.vector-squares}
$v^{2} = Q(v)$ and $uv + vu = 2B(u, v)$, so $uv = B(u, v) + \tfrac12(uv - vu)$.
\emph{Source:} the algebra core of the Transport Lab, proved there with an exact check; standard, see C. Doran and A. Lasenby, Geometric Algebra for Physicists, Cambridge University Press (2003), Ch. 4 (foundations of geometric algebra).
\end{fact}

\begin{fact}[The group of null rotations]\label{c:transport.group-null-rotations}
The null rotations $R_u = 1 + vu$ with $u \in U$ form an abelian subgroup of $\Spin^0(p,q)$ that fixes $v$ and is isomorphic to $(U, +)$, equivariantly for the rotors of $U$.
\emph{Source:} T1, Lem 4.6 (\href{https://doi.org/10.5281/zenodo.21743863}{doi:10.5281/zenodo.21743863}).
\end{fact}

\begin{fact}[Totality]\label{c:transport.totality}
On the quadric of the completion, $\iota[N] = [e_+ N e_+]$ is a total continuous involution that equals the reciprocal $x \mapsto x/Q(x)$ on invertible vectors and sends each $x_0$ with $Q(x_0) = 0$ to $[x_0 + \frac12 \nInf]$, and it is the only such continuous extension.
\emph{Source:} T1, Thm 4.10 (\href{https://doi.org/10.5281/zenodo.21743863}{doi:10.5281/zenodo.21743863}).
\end{fact}

\section*{Conclusion}
Starting from points as null vectors, the trail has recovered the conformal spinor space, the Klein correspondence, the incidence relation, the reality condition and local spinor transport, all from the completion itself.

No postulate is needed on this stretch; the one standing condition is that the conformal manifold be spin, which the transport step assumes. The vacuum-line picture of spacetime that emerges is the ground the later subjects build on.

\end{document}